% Paper insert: no documentclass, document environment, or bibliography.
% Load bittide_math_preamble.tex in the host preamble.
% Merge bittide_sources.bib into the host bibliography.
\section{Proposed Fuzzy PI Control for bittide and Nonlinear Analysis}
\label{sec:fpi_controller}

We propose a decentralized, occupancy-based incremental Fuzzy PI controller
for bittide. Each node combines receiving-buffer error with its
update-to-update change to compute an oscillator-correction increment.
We specify the complete controller and its continuous-clock interface,
establish a comparable PI baseline and derive the nonlinear properties
used in the reference-model analysis.

\subsection{Continuous-Clock Interface and Proposed Fuzzy PI}
\label{subsec:fpi_laws}

The plant remains the continuous-phase Callisto model within the bittide
framework~\cite{lall2024logicalSynchrony}. With constant
intrinsic frequency $\omega_i^u$, physical oscillator phase $\theta_i(t)$
evolves between actuation events according to
\begin{align}
 \dot\theta_i(t)&=\omega_i(t)=\omega_i^u+u_i(t),\nonumber\\
 u_i(t)&=u_{i,k},\qquad t\in[t^a_{i,k},t^a_{i,k+1}).
 \label{eq:fpi_continuous_plant}
\end{align}
Positive physical frequencies are an admissibility requirement.
The superscript $a$ denotes actuation. Phase is continuous; a controller
update changes its slope, not its value. The local polling period $p>0$
and execution delay $d_{\rm ctrl}\ge0$ are measured in local-clock phase
units. If $t^m_{i,k}$ denotes the measurement time used at actuation $k$,
the native scheduling relations are
\begin{align}
 \theta_i(t^a_{i,k+1})-\theta_i(t^a_{i,k})&=p,\nonumber\\
 \theta_i(t^a_{i,k})-\theta_i(t^m_{i,k})&=d_{\rm ctrl}.
 \label{eq:fpi_local_schedule}
\end{align}
The first measurement is at $t^m_{i,0}=0$; with zero initial correction,
the first actuation is $t^a_{i,0}=d_{\rm ctrl}/\omega_i^u$.
Before it, $u_i(t)=0$. Thereafter the physical interval is
$p/(\omega_i^u+u_{i,k})$ in the constant-intrinsic-frequency case, so
polling is not a common fixed physical-time grid.

Frames are counted at integer phase crossings. In normal operation, with
the input/credit guards open, the occupancy obeys the count balance
\begin{equation}
 \beta_{j\to i}(t)=\beta_{j\to i}(0)+A^{\rm arr}_{ji}(t)-N_i(t),
 \label{eq:fpi_count_balance}
\end{equation}
where $A^{\rm arr}_{ji}(t)$ counts arrivals during $(0,t]$, including
initial in-flight frames, and
$N_i(t)=\lfloor\theta_i(t)\rfloor-\lfloor\theta_i(0)\rfloor$ counts
local consumption. Arrival ordering and initialization retain the Callisto
marking and prehistory. Its variable-latency extension preserves frame
order and may release batches; it assumes no finite serialization rate.
For a constant physical path delay $\tau_{ji}\ge0$, the arrival count is
\begin{equation}
 A^{\rm arr}_{ji}(t)=\lfloor\theta_j(t-\tau_{ji})\rfloor
                     -\lfloor\theta_j(-\tau_{ji})\rfloor,
 \label{eq:fpi_fixed_delay_arrivals}
\end{equation}
using the declared negative-time phase history.
If a guard blocks consumption, actual data-progress counts replace $N_i$;
the physical oscillator continues. Thus discrete frame events and sampled
controller updates are embedded in continuous physical time. Fluid
occupancies and a common sampling grid are introduced only for Section~\ref{sec:btd_reference}'s
reference proof.

Let $\mathcal N_i^-$ denote node $i$'s incoming-neighbour set. We consider fixed, nonempty incoming neighbour sets with all link-status
masks active. A changing link mask would require consistent aggregation in
both controllers. Let $k$ index the successive measurements processed by node $i$, with measured
incoming occupancies $\beta_{ji,k}=\beta_{j\to i}(t^m_{i,k})$ and fixed targets $\beta^*_{ji}$.
The subscript $ji$ abbreviates receiving edge $j\to i$; $k$ is a node-local
measurement index, not a shared physical sampling time. Define
the sum error, measured in frames, by
\begin{equation}
 E_{i,k}=\sum_{j\in\mathcal N_i^-}(\beta_{ji,k}-\beta^*_{ji}).
 \label{eq:fpi_error}
\end{equation}
An excess occupancy gives positive $E_{i,k}$. This frame error is used to
calibrate the PI comparison; the proposed controller operates on the
following percentage-point error.

For incoming degree $d_i=|\mathcal N_i^-|>0$ and common capacity $C$, let
\begin{equation*}
\begin{aligned}
 \bar p_{i,k}&=d_i^{-1}\sum_{j\in\mathcal N_i^-}100\beta_{ji,k}/C,\\
 r_i&=d_i^{-1}\sum_{j\in\mathcal N_i^-}100\beta^*_{ji}/C.
\end{aligned}
\end{equation*}
Here $r_i$ is the setpoint. The percentage-point error is
\begin{equation}
 e_{i,k}=r_i-\bar p_{i,k}=-\gamma_iE_{i,k},
 \qquad \gamma_i=\frac{100}{d_iC}.
 \label{eq:fpi_percent}
\end{equation}
The experiments with midpoint targets use $r_i=50$. The same identity also
holds for a target-centered percentage interface. Unequal inbound capacities
instead produce a weighted token sum, so~\eqref{eq:fpi_percent} must not be
assumed for a heterogeneous percentage average.

The incremental Fuzzy controller computes
\begin{align}
 X_{i,k}&=K_e e_{i,k},\quad
 Y_{i,k}=K_{\Delta e}(e_{i,k}-e_{i,k-1}),\nonumber\\
 A_{i,k}&=A_{i,k-1}+K_u F(X_{i,k},Y_{i,k}).
 \label{eq:fpi_fuzzy_law}
\end{align}
The requested correction is $q^{\rm F}_{i,k}=A_{i,k}$.
Thus $Y$ uses a sample difference, not a time-normalized derivative.
We use $K_e,K_{\Delta e}>0$ and $K_u<0$. The state $A$ is the total
requested frequency correction; $K_u$ converts a normalized surface
output into a frequency-correction increment.

Both inputs use negative large (NL), negative small (NS), zero (ZE), positive small (PS) and positive large (PL) labels, with breakpoints
$-10,-5,0,5,10$, triangular middle sets and shoulder outer sets.
For $[z]_+=\max\{z,0\}$ and
$\operatorname{clip}_{[0,1]}z=\min\{1,\max\{0,z\}\}$, the full
membership functions are
\begin{align}
 \varphi_{\rm NL}(z)&=\operatorname{clip}_{[0,1]}((-z-5)/5),\nonumber\\
 \varphi_{\rm NS}(z)&=[1-|z+5|/5]_+,\nonumber\\
 \varphi_{\rm ZE}(z)&=[1-|z|/5]_+,\nonumber\\
 \varphi_{\rm PS}(z)&=[1-|z-5|/5]_+,\nonumber\\
 \varphi_{\rm PL}(z)&=\operatorname{clip}_{[0,1]}((z-5)/5).
 \label{eq:fpi_full_memberships}
\end{align}
Each membership is an activation degree in $[0,1]$; a singleton is a
numerical rule consequent. Table~\ref{tab:fpi_full_rules} specifies all
25 consequents. For label indices $a,b\in\{1,\ldots,5\}$, the inference is
\begin{align}
 w_{ab}(X,Y)&=\min\{\varphi_a(X),\varphi_b(Y)\},\nonumber\\
 \ell(a,b)&=\min\{5,\max\{1,a+b-3\}\},\nonumber\\
 F(X,Y)&=\frac{\sum_{a,b}w_{ab}(X,Y)s_{\ell(a,b)}}
                   {\sum_{a,b}w_{ab}(X,Y)},
 \label{eq:fpi_rules}
\end{align}
where $(s_1,\ldots,s_5)=(-5,-2.5,0,2.5,5)$. Rules with the same
consequent remain separate terms in both sums; there is no intermediate
maximum aggregation. This normalized weighted average implies $|F|\le5$,
but does not by itself bound the accumulated state $A$.

\begin{table}[t]
\caption{Rule consequents of the proposed bittide Fuzzy PI controller.}
\label{tab:fpi_full_rules}
\centering
\begin{tabular}{@{}c|ccccc@{}}\toprule
$X\backslash Y$ & NL & NS & ZE & PS & PL\\\midrule
NL & NL & NL & NL & NS & ZE\\
NS & NL & NL & NS & ZE & PS\\
ZE & NL & NS & ZE & PS & PL\\
PS & NS & ZE & PS & PL & PL\\
PL & ZE & PS & PL & PL & PL\\\bottomrule
\end{tabular}
\end{table}

Applied corrections are obtained through the common execution layer, for
example $u_{i,k}=\operatorname{sat}_{u_{\max}}(q_{i,k})$, where $u_{\max}>0$ and $\operatorname{sat}_{b}(z)=\min\{b,\max\{-b,z\}\}$. The physical frequency is $\omega_i^u+u_i$, with constant intrinsic frequency $\omega_i^u$.
After each ordinary inference, the controller stores $e_{i,k}$ as its
previous error. Initialize $A_{i,-1}=0$ and $e_{i,-1}=0$. The reported
normal runs retain the raw accumulator if external clipping occurs
(anti-windup is disabled); they do not silently overwrite it with the
limited command. The update order is therefore: process the local
measurement, form $e,X,Y$, evaluate the memberships and rules, update
$A$, store the previous error, limit the request and apply the held
correction at $t^a_{i,k}$. The memberships, consequents and scaling gains
remain fixed during a run; the correction changes with the measured
inputs. No observer or FFP mode is required by this law.

\subsection{PI Baseline, Scale Calibration and the Linear Counterpart}
\label{subsec:fpi_calibration}

The occupancy-feedback PI baseline follows the bittide control setting
of~\cite{lall2022resistance}. With positive gains $k_p,k_i$, the native
implementation integrates before output:
\begin{align}
 I_{i,k}&=I_{i,k-1}+E_{i,k},\nonumber\\
 q^{\rm PI}_{i,k}&=k_pE_{i,k}+\kappa_I I_{i,k},\quad\kappa_I=k_i h.
 \label{eq:fpi_native}
\end{align}
Here $I$ is the PI sample sum and $q$ the requested correction. The fixed
integration scale is $h=p/\omega_{\rm base}$, where
$\omega_{\rm base}$ is the implementation reference frequency.
It is not each actual physical sampling interval. Defining
$\Delta q_{i,k}=q_{i,k}-q_{i,k-1}$ gives
\begin{equation}
 \Delta q^{\rm PI}_{i,k}
 =k_p(E_{i,k}-E_{i,k-1})+\kappa_I E_{i,k}.
 \label{eq:fpi_pi_increment}
\end{equation}
Define the globally linear counterpart
\begin{equation}
 F_0(X,Y)=\frac{X+Y}{2}.
 \label{eq:fpi_f0}
\end{equation}
Let $\Delta A^0_{i,k}=A^0_{i,k}-A^0_{i,k-1}$. Replacing $F$ by $F_0$ in~\eqref{eq:fpi_fuzzy_law} gives
\begin{align}
 \Delta A^{0}_{i,k}
 &=-\frac{K_u\gamma_i}{2}\bigl[K_eE_{i,k}
                  +K_{\Delta e}(E_{i,k}-E_{i,k-1})\bigr].
 \label{eq:fpi_substitution}
\end{align}
Consequently the exact coefficient correspondence is
\begin{equation}
 \boxed{k_p=-\frac{K_u\gamma_iK_{\Delta e}}2,\qquad
        \kappa_I=-\frac{K_u\gamma_iK_e}2.}
 \label{eq:fpi_matching}
\end{equation}
For prescribed positive $k_p,k_i,h$ and a chosen $K_e>0$, one may use
\begin{equation}
 K_u=-\frac{2k_ih}{\gamma_iK_e},\qquad
 K_{\Delta e}=\frac{K_ek_p}{k_ih}.
 \label{eq:fpi_design_gains}
\end{equation}
Changing capacity, degree or $h$ changes this calibration. Uniform sum-PI
gains can require node-specific Fuzzy parameters on a nonuniform network.

\begin{btdproposition}[Event-level linear-counterpart equivalence]
\label{prop:fpi_equivalence}
Suppose~\eqref{eq:fpi_percent} and~\eqref{eq:fpi_matching} hold. For an
actual or fictitious previous error $E_{i,-1}$, initialize
\begin{equation}
 A^0_{i,-1}=k_pE_{i,-1}+\kappa_I I_{i,-1},\qquad
 e_{i,-1}=-\gamma_iE_{i,-1}.
 \label{eq:fpi_initialization}
\end{equation}
Then the $F_0$ controller and native PI produce identical requested
corrections for every common measurement sequence.
\end{btdproposition}
\begin{proof}
Substitution~\eqref{eq:fpi_substitution} and calibration give exactly
\eqref{eq:fpi_pi_increment}. The initial relation in
\eqref{eq:fpi_initialization} and induction preserve
$A^0_{i,k}=k_pE_{i,k}+\kappa_I I_{i,k}$ at every update.
\end{proof}

Zero constructor states correspond to the fictitious $E_{i,-1}=0$.
Even when $E_{i,0}\ne0$, both first outputs equal
$(k_p+\kappa_I)E_{i,0}$. Initializing only the Fuzzy previous error to
$e_{i,0}$ would remove its first proportional contribution and break the
comparison. In a deterministic causal execution, identical plant history,
external traces, sampling/actuation rules, tie-breaking and compatible
limiting/reset policies extend the identity to the closed-loop event
sequence, up to numerical arithmetic. Physical sampling need not be globally
synchronous. $F_0$ must remain unclipped internally; limiting it to the
singleton range would destroy global PI equivalence. Axis calibration
provides a common reference, rather than making full Fuzzy PI linear or
establishing an optimally tuned PI baseline.

\subsection{Exact Nonlinear Mechanism}
\label{subsec:fpi_mechanism}

For $0<r\le5$, define the central diamond and auxiliary quantities
\begin{align}
 \Omega_r&=\{(X,Y):|X|+|Y|\le r\},\nonumber\\
 m&=\min\{|X|,|Y|\},\quad s=|X|+|Y|.
 \label{eq:fpi_domain}
\end{align}
Let the same-sign interaction be
\begin{equation}
 J(X,Y)=\begin{cases}
 \operatorname{sgn}(X)m,&XY>0,\\
 0,&XY\le0.
 \end{cases}
 \label{eq:fpi_interaction}
\end{equation}

\begin{btdproposition}[Exact central surface]
\label{prop:fpi_central}
For every $(X,Y)\in\Omega_5$, the implemented rule law satisfies
\begin{equation}
 \boxed{F(X,Y)=\frac{F_0(X,Y)+J(X,Y)}{1+2m/5}.}
 \label{eq:fpi_central}
\end{equation}
\end{btdproposition}
\begin{proof}
Put $a=|X|/5$ and $b=|Y|/5$. Since $a+b\le1$, only the central and
adjacent sign-labelled sets are active on each input. Their four rule
weights are $1-\max(a,b)$, $a$, $b$ and $\min(a,b)$, which sum to
$1+2\min(a,b)$. The single-nonzero-input consequents contribute
$X/2+Y/2$. The joint consequent contributes
$\operatorname{sgn}(X)m$ for equal signs and zero for opposite signs.
Dividing the summed contributions proves the formula, including boundaries
by continuity.
\end{proof}

For same-sign inputs the surface-to-reference ratio is
\begin{equation}
 \frac{F}{F_0}=\frac{1+2m/s}{1+2m/5}\in[1,2],\qquad XY>0.
 \label{eq:fpi_same_sign}
\end{equation}
Here $s\le5$ gives the lower bound and $m\le s/2$ gives the upper bound.
For opposite signs,
\begin{equation}
 F=\frac{F_0}{1+2m/5},\qquad XY<0.
 \label{eq:fpi_opposite_sign}
\end{equation}
The surface therefore enhances a same-sign increment and attenuates the
magnitude of an opposing-sign increment relative to the same-input $F_0$.
Axes give $F=F_0$; when $X+Y=0$, both outputs vanish. Since $K_u<0$,
positive $F$ produces a negative correction increment.

\begin{btdcorollary}[Regional pointwise sector]
\label{cor:fpi_sector}
On $\Omega_r$, the surface admits $F=\alpha(X,Y)F_0$ with
\begin{equation}
 \frac1{1+r/5}\le\alpha(X,Y)\le2.
 \label{eq:fpi_sector}
\end{equation}
Equivalently, including $F_0=0$ without division,
\begin{equation}
 \left(F-\frac{F_0}{1+r/5}\right)(2F_0-F)\ge0.
 \label{eq:fpi_sector_product}
\end{equation}
\end{btdcorollary}
\begin{proof}
Use~\eqref{eq:fpi_same_sign} for equal signs and
$m\le r/2$ in~\eqref{eq:fpi_opposite_sign} for opposite signs.
At $F_0=0$, the exact formula gives $F=0$ and any admissible $\alpha$
can be selected.
\end{proof}

This is a state-dependent pointwise representation, not a con\-stant-gain
linearization. In particular, $F(X,0)=X/2$ and $F(0,Y)=Y/2$, whereas
$F(t,t)/t\to2$ as $t\downarrow0$; the origin has no ordinary Jacobian.
For identical scaled input pairs, including consistent previous-error
initialization, and a common initial accumulator, define
$\mathcal R_F=F-F_0$. With no reset or accumulator projection,
\begin{equation}
 A^{\rm F}_{i,k}-A^0_{i,k}
 =K_u\sum_{\ell=0}^{k}\mathcal R_F(X_{i,\ell},Y_{i,\ell}).
 \label{eq:fpi_memory}
\end{equation}
Thus a few nonlinear increments can leave persistent correction memory.
Equation~\eqref{eq:fpi_memory} is a same-input accounting identity; separate
closed loops generally generate different measurements and schedules.

The mechanism enables direction-dependent response shaping at calibrated
scales. Its pointwise properties alone do not guarantee lower jitter,
smaller buffer peaks or stability. Those questions require an invariant
closed-loop domain and a model of the plant and execution timing, as
developed in the reference analysis and evaluated experimentally.
